\documentclass[11pt]{article}
\usepackage{fontspec}
\usepackage[french,english]{babel}
\usepackage{amsmath}
\usepackage{amssymb}
\usepackage{amsthm}
\usepackage{mathrsfs}
\usepackage{tikz-cd}
\usepackage{enumitem}

\newcommand{\printedpage}[1]{\marginpar{\footnotesize p.~#1}}
\newcommand{\msp}[1]{\marginpar{\footnotesize [#1]}}
\newcommand{\Fl}{\mathrm{Fl}}
\newcommand{\M}{\mathcal{M}}

\begin{document}

\printedpage{--} % no page number printed on this page
% head: (none)

\begin{center}
\vspace*{5cm}
{\LARGE\sffamily THE DERIVATORS}

\vspace{1.5cm}
{\large\textbf{ALEXANDRE GROTHENDIECK}}

\vspace{2.5cm}
{\large\textbf{Chapter XIII}}

\vspace{0.5cm}
{\large\textbf{Categories of models (1)}}
\end{center}

\clearpage

\printedpage{--} % no page number printed on this page
% head: (none)

\vspace*{6cm}
\noindent This text has been deciphered and transcribed in \LaTeXe{} by M. Künzer. It has been edited by M. Künzer, J. Malgoire, and G. Maltsiniotis. The transcription is as faithful as possible to the manuscript. For the few obvious corrections, or rare comments by the editors, as well as for the original numbering of the manuscript pages, typewriter characters \texttt{[typewriter]} in brackets are used. A question mark in brackets means that one is not sure of the preceding word. The indexes, the list of axioms, the glossary, and the bibliographical references at the end of the chapter have been compiled by the editors. The page numbers in the indexes are those of the original.

\noindent This edition is provisional. Remarks, comments and corrections are welcome. Send a message to:

\noindent\textbf{[e-mail address removed]}

\hfill G. Maltsiniotis\hspace*{2cm}

\clearpage

\printedpage{1}
% head: A. Grothendieck : Dérivateurs, Ch. XIII, Édition des Universités de Montpellier II et Paris VII
% foot: Version 25 janvier 2004

\msp{page 1}

\section*{1\quad Quillen couples $(\Phi, \Psi)$}

I want first to give a weakened variant of the notion of a Quillen triple (Quillen's ``model category''). This will not be an entirely self-dual notion -- thus there will be \textit{two} weakened variants, dual to one another.

But even before that, a sorites on the LLP (left lifting property) and RLP (right lifting property) (\footnote[1]{give formal definition, with reference to \textsc{Quillen}.}). Let $\M$ be any category (most often, $\M$ will be large). We consider subsets
\[
\Phi,\ \Psi\ \subset\ \Fl(\M)
\]
and we define maps
\[
\boxed{\begin{array}{ccc}
\Phi & \longmapsto & \Phi_* \\[1ex]
\Psi & \longmapsto & \Psi^*
\end{array}}
\qquad
\Fl(\M) \longrightarrow \Fl(\M)
\]
\texttt{[rather} $\mathfrak{P}(\Fl(\M) \longrightarrow \mathfrak{P}(\Fl(\M))$\texttt{,}\newline\texttt{where for every set} $E$\texttt{,} $\mathfrak{P}(E)$ \texttt{denotes the set of subsets of} $E$\,\texttt{]} by % [sic] \mathfrak{P}(\Fl(\M) -> \mathfrak{P}(\Fl(\M)), a closing parenthesis is missing in the first term
\[
\begin{aligned}
\Phi_* \ &= \ \{\text{set of } p \in \Fl(\M) \text{ that have the RLP with respect to } \Phi\} \\
\Psi^* \ &= \ \{\text{set of } i \in \Fl(\M) \text{ that have the LLP with respect to } \Psi\},
\end{aligned}
\]
and then
\[
\boxed{\begin{array}{ccc}
\Phi & \longmapsto & \widetilde{\Phi} \\[1ex]
\Psi & \longmapsto & \overline{\Psi}
\end{array}}
\qquad
\Fl(\M) \longrightarrow \Fl(\M)
\]
\texttt{[rather} $\mathfrak{P}(\Fl(\M) \longrightarrow \mathfrak{P}(\Fl(\M))$\texttt{]} by % [sic] \mathfrak{P}(\Fl(\M) -> \mathfrak{P}(\Fl(\M)), a closing parenthesis is missing in the first term
\[
\begin{aligned}
\widetilde{\Phi} \ &= \ (\Phi_*)^* \\
\overline{\Psi} \ &= \ (\Psi^*)_* \ .
\end{aligned}
\]

\msp{page 2}

I summarize the most immediate properties of these operations in a

\noindent\textbf{Proposition 1.}

\begin{enumerate}[label={}]
\item[a)] $\Phi \longmapsto \Phi_*$, $\Psi \longmapsto \Psi^*$ \textit{are decreasing maps from} $\Fl(\M)$ \texttt{[rather} $\mathfrak{P}(\Fl(\M))$\texttt{]} \textit{to itself, for the order relation} $\subset$\textit{, and} $\Phi \longmapsto \widetilde{\Phi}$\textit{,} $\Psi \longmapsto \overline{\Psi}$ \textit{increasing maps.}

\printedpage{2}
% head: A. Grothendieck : Dérivateurs, Ch. XIII, Édition des Universités de Montpellier II et Paris VII
% foot: Version 25 janvier 2004

\item[b)] \textit{We have}
\[
\Phi\ \subset\ \widetilde{\Phi},
\]
\textit{and for $\Phi = \widetilde{\Phi}$ to hold, it is necessary and sufficient that $\Phi$ be of the form $\Psi^*$. We then say that $\Phi$ is} \textup{left closed} (\footnote[2]{Give formal definition.})\textit{. The set of subsets of} $\Fl(\M)$ \textit{that are left closed is stable under arbitrary intersections} (\footnote[3]{more precisely, $\bigcap_i(\Psi_i^*) = (\bigcup_i \Psi_i)^*$}), \textit{and $\widetilde{\Phi}$ is the smallest closed subset of} $\Fl(\M)$ \textit{that contains $\Phi$ and the largest subset $\Phi'$ of} $\Fl(\M)$ \textit{such that $\Phi'_* = \Phi_*$ (or such that $\Phi'_* \supset \Phi_*$).}

\textit{Symmetric statements for $\Psi$: we have}
\[
\Psi\ \subset\ \overline{\Psi},
\]
\textit{we have $\Psi = \overline{\Psi}$ if and only if $\Psi$} \texttt{[is]} \textit{of the form $\Phi_*$, we then say that $\Psi$ is} \textup{right closed.} \textit{The set of subsets of} $\Fl(\M)$ \textit{that are right closed is stable under arbitrary intersections} (\footnote[4]{more precisely, we have $\bigcap_i(\Phi_{i\,*}) = (\bigcup_i \Phi_i)_*$}), \textit{and $\overline{\Psi}$ is the smallest right closed subset containing $\Psi$, and also the largest subset $\Psi'$ of} $\Fl(\M)$ \textit{such that $\Psi'^* = \Psi^*$ (or such that $\Psi'^* \supset \Psi^*$).}

\item[] \hspace*{-2.5em}\msp{page 3}

\item[c)] \textit{The maps}
\[
\Phi \ \longmapsto\ \Phi_*, \qquad\quad \Psi\ \longmapsto\ \Psi^*
\]
\textit{establish mutually inverse bijections between the set of left closed subsets $\Phi$ of} $\Fl(\M)$\textit{, and the set of right closed subsets $\Psi$. This is an anti-isomorphism of ordered sets. We then have}
\[
(\Phi \cap \Phi')_* \ = \quad
\begin{array}{l}
\textit{right closed subset of } \Fl(\M) \ \textit{``generated by } \Phi_*,\ \Phi'_*\textit{'',} \\
\text{i.e. \texttt{[the]} \textit{smallest right closed subset containing} } \Phi_*, \\
\Phi'_*, \text{ i.e. } \overline{\Phi_* \cup \Phi'_*}
\end{array}
\]
\texttt{[for} $\Phi$\texttt{,} $\Phi'$ \texttt{left closed subsets]}\textit{, more generally}
\[
(\bigcap_i \Phi_i)_* \ = \ \overline{\bigcup_i (\Phi_{i\,*})}
\]
\textit{for every family $(\Phi_i)_{i \in I}$ of left closed subsets of} $\Fl(\M)$\textit{. Dual formulas} \texttt{[for right closed subsets]}
\[
\left\{
\begin{aligned}
(\Psi \cap \Psi')^* \ &= \ \widetilde{\Psi^* \cup \Psi'^*} \\
(\bigcap_j \Psi_j)^* \ &= \ \widetilde{\bigcup_j (\Psi_j^*)}\ .
\end{aligned}
\right.
\]

\item[d)] \textit{If} $\Phi, \Psi \subset \Fl(\M)$\textit{, then we have}
\[
\Phi \subset \Psi^* \iff \Psi \subset \Phi_*, \qquad\quad \textit{we then write } \Phi \longleftrightarrow \Psi,
\]
\textit{this means that $\Phi$ has the LLP for $\Psi$, or that $\Psi$ has the RLP for $\Phi$.}
\end{enumerate}

\printedpage{3}
% head: A. Grothendieck : Dérivateurs, Ch. XIII, Édition des Universités de Montpellier II et Paris VII
% foot: Version 25 janvier 2004

\msp{page 4}

\noindent\textbf{Proposition 2.} \textit{Let $\Phi$ be a left closed subset of} $\Fl(\M)$\textit{. Then $\Phi$ satisfies the conditions} a) b) c) \textit{below:}

\begin{enumerate}[label={}]
\item[a)] \textit{$\Phi$ contains the isomorphisms, and is stable under composition, and under cobase change (}i.e. \textit{if we have}
$\vcenter{\hbox{\begin{tikzcd}
A \arrow[r, "i"] \arrow[d] & B \arrow[d] \\
A' \arrow[r, "i'"] & B'
\end{tikzcd}}}$
\textit{a cocartesian square in $\M$, with $i \in \Phi$, we have $i' \in \Phi$).}

\item[b)] \textit{$\Phi$ stable under direct factors (as a subset of the set of objects of} $\underline{\Fl}(\M)$\textit{).}

\item[c)] \textit{$\Phi$ ``stable under ordinal direct limits'', in the following sense: Let $I$ be a well-ordered set, $(A_i)_{i \in I}$ a direct system in $\M$, such that the $u_{ji} : A_i \longrightarrow A_j$ are in $\Phi$, and $A_{i_0} \xleftarrow{\ \sim\ } \varinjlim\limits_{i<i_0} A_i$ if $i_0 \in I$ is a limit element. Suppose that $A = \varinjlim\limits_{I} A_i$ exists. Then $\forall\ i_0 \in I$, the canonical arrow $A_{i_0} \longrightarrow A$ is in $\Phi$.}
\end{enumerate}

There is a dual statement for a right closed subset $\Psi$; it contains the isomorphisms, is stable under composition and under base change, under direct factors, and finally stable under ordinal inverse limits.

\msp{page 5}

We will have a partial converse in accessible categories. It is for the statement of this converse that I have set apart condition \textit{b}) of stability under direct factors, instead of lumping it together with \textit{a}).

\noindent\textbf{Definition 1.}

\begin{enumerate}[label={}]
\item[\textit{a})] (For the record.) RLP, LLP, the relation $\Phi \longleftrightarrow \Psi$ between subsets of $\Fl(\M)$ (nonsymmetric relation!).

\item[\textit{b})] (For the record.) Subset of $\Fl(\M)$ left closed, right closed.

\item[\textit{c})] Let $\Phi \longleftrightarrow \Psi$. We say that the pair $(\Phi, \Psi)$ has the \textit{Quillen factorization property}, if \textit{every} arrow $f$ of $\M$ factors as
\[
f = pi, \qquad\quad i \in \Phi,\ p \in \Psi.
\]
We say that $(\Phi, \Psi)$ is a \textit{Quillen couple} if it is a pair associated by \texttt{[the]} properties LLP and RLP ($\Phi \longleftrightarrow \Psi$), and \texttt{[if]} they \texttt{[the subsets} $\Phi$ \texttt{and} $\Psi$\texttt{]} are closed (\textit{i.e.} $\Phi = \Psi^*$, $\Psi = \Phi_*$), and if they have the Quillen factorization property.
\end{enumerate}

I state right now the essential technical result, ensuring the existence of many Quillen couples:

\end{document}
